\begin{table}[!htbp]
\centering
\caption{Robustness to regional spillovers, 9th LP: ATT by relative period}
\label{tab:ap_event_spillover_9th_lp}
\vspace{-0.15cm}
\scriptsize
\setlength{\tabcolsep}{4pt}
\renewcommand{\arraystretch}{0.92}
\begin{tabular}{lccc}
\toprule
Relative period & Main & Exclude exposed DE & Exclude exposed district-DE \\
\midrule
-4 & 0.019 (0.060) & 0.048 (0.058) & 0.041 (0.057) \\
-3 & 0.016 (0.043) & -0.036 (0.055) & -0.066 (0.049) \\
-2 & 0.003 (0.030) & -0.065 (0.041) & -0.053 (0.039) \\
-1 & 0.000 & -0.013 (0.033) & -0.002 (0.031) \\
0 & 0.303 (0.039) & 0.310 (0.042) & 0.324 (0.042) \\
1 & 0.425 (0.038) & 0.426 (0.045) & 0.421 (0.039) \\
2 & 0.528 (0.040) & 0.495 (0.064) & 0.538 (0.041) \\
3 & 0.598 (0.046) & 0.546 (0.079) & 0.609 (0.049) \\
4 & 0.698 (0.051) & 0.631 (0.066) & 0.737 (0.054) \\
\bottomrule
\end{tabular}
\begin{flushleft}
\footnotesize Note: each cell reports the ATT and standard error in parentheses, in standard deviations of annual proficiency. Period \(k=-1\) is normalized as the baseline and therefore appears without a standard error.
\end{flushleft}
\end{table}



